Accurate channel state information (CSI) in OFDM links is obtained by transmitting pilot subcarriers and estimating the channel from the resulting observations, but every pilot consumes time--frequency resources and adds latency. Conventional pilot designs use a fixed, pre-determined subcarrier set and then run a batch estimator, such as the linear minimum mean-square error (LMMSE) method, committing to a pilot budget without adapting to information already gained. We propose an incremental pilot-selection framework that measures one subcarrier at a time, updates an LMMSE posterior after each measurement, and selects the next subcarrier via a model-based active rule that greedily maximizes the predicted reduction in the posterior error-covariance trace. Experiments under matched and mismatched channel priors, including correlated Gaussian and 3GPP tapped-delay-line channels, show that active selection reaches a target estimation quality with fewer pilots than fixed or random selection. These results also expose a central practical challenge: the active rule and its stopping criterion both rely on an assumed channel covariance, which is rarely known exactly in practice and must instead be estimated or approximated.
